\documentclass[a4paper,11pt]{article}
\usepackage{cmap}
\usepackage[T2A]{fontenc}
\usepackage[utf8]{inputenc}
\usepackage[russian]{babel}
\usepackage[margin=22mm]{geometry}
\usepackage[hidelinks]{hyperref}
\hypersetup{pdftitle={Задача 14. Две метки}}
\setlength{\parindent}{0pt}
\setlength{\parskip}{5pt}

\begin{document}

\begin{center}
  {\large\bfseries Задача 14. Две метки}
\end{center}

\textbf{Условие.} Двух друзей взяли в плен. Каждому поставят на лбу метку,
цвет которой выберут случайно: чёрный или белый. Каждый сможет увидеть
только метку другого. Затем друзья одновременно назовут по одному цвету.
Если хотя бы один из них угадает цвет собственной метки, обоих отпустят.

До того как им поставят метки, друзья могут договориться о стратегии.
После этого общаться будет нельзя. Как им договориться, чтобы обоих
наверняка отпустили?

\smallskip
{\small Источник: u/ConnectingTheDotsCh,
\href{https://www.reddit.com/r/puzzles/comments/1y4jw6m/guess_the_color_of_the_mark/}{\emph{Guess the color of the mark!}} (Reddit).}

\clearpage

\textbf{Решение.} Пусть первый называет тот цвет, который видит у второго.
Если метки одного цвета, первый угадает. В каком случае всё зависит
от ответа второго? Когда первый ошибается, то есть когда метки разного цвета.

Значит, второму достаточно всегда называть цвет, противоположный тому,
который он видит у первого. Тогда при одинаковых метках угадает первый,
а при разных~--- второй. В любом случае обоих отпустят.

\end{document}
